\chapter{Problem Formulation}
\label{sec:problem}

This chapter formulates the dynamics modeling problem considered in this work. The formulation captures environmental variability, specifies the available observations, and defines a trajectory-level predictive learning objective.

\section{System Dynamics}

Consider a discrete-time controlled dynamical system. Let $t \in \mathbb{N}$ index discrete time steps, $s_t \in \mathcal{S} \subset \mathbb{R}^{n_s}$ denote the system state, and $a_t \in \mathcal{A} \subset \mathbb{R}^{n_a}$ denote the control input.

In off-road driving scenarios, system dynamics depend not only on the current state and control input, but also on environmental factors such as terrain properties and tire--ground interaction. These factors are not directly observable and are modeled as a latent variable $c_t \in \mathcal{C}$.

The system evolves according to:
\begin{equation}
  s_{t+1} = f^*(s_t, a_t;\, c_t) + w_t, \qquad w_t \sim \mathcal{W},
  \label{eq:true-dynamics}
\end{equation}
where $f^*: \mathcal{S} \times \mathcal{A} \times \mathcal{C} \to \mathcal{S}$ is the unknown dynamics function and $w_t$ represents process noise.

Because the environmental variable $c_t$ is not directly measurable, its effect on system dynamics must be inferred indirectly. The learned models considered in this work are deterministic approximations of the underlying stochastic dynamics.

\section{Observations and Dataset}

In addition to state and control measurements, auxiliary observations are available, which provide indirect information about the latent environmental variable $c_t$. These observations may include high-dimensional sensory inputs.

An offline dataset of trajectory segments is assumed:
\begin{equation}
  \mathcal{D}_{\text{train}} =
  \left\{
  \big(s_{t:t+H}^{(i)},\, a_{t:t+H-1}^{(i)},\, \xi^{(i)}\big)
  \right\}_{i=1}^{N},
\end{equation}
where $H$ denotes the prediction horizon, and $\xi^{(i)}$ denotes context features associated with the beginning of the $i$-th trajectory segment.

Each tuple corresponds to a trajectory segment collected under a particular environmental condition. The horizon $H$ is chosen to match the rollout length used during downstream planning.

The dataset is collected from the physical system under varying environmental conditions and remains fixed for all models. Auxiliary observations are available for models that condition parameters on environmental context, but are not required by fixed-parameter baselines.

\section{Context-Conditioned Dynamics Modeling}

The objective is to learn a predictive model of the system dynamics that accounts for environmental variability.

Models of the following form are considered:
\begin{equation}
  \hat{s}_{t+1} = \hat{f}(s_t, a_t;\, \phi^{(i)}),
  \label{eq:learned-model}
\end{equation}
where $\hat{f}$ denotes a parameterized dynamics model whose structure is specified a priori, and $\phi^{(i)}$ denotes model parameters associated with the $i$-th trajectory segment.

The formulation assumes a Markovian representation, while allowing practical implementations to maintain additional internal state.

The parameter $\phi^{(i)}$ may be either constant or context-dependent. Context-dependent parameterizations allow the model to adapt to changes in the underlying environmental conditions.

To incorporate contextual information, $\phi^{(i)}$ is modeled as a function of context features:
\begin{equation}
  \phi^{(i)} = \psi(\xi^{(i)}),
\end{equation}
where $\psi$ is a mapping from context features to model parameters.

For each trajectory segment, the parameter vector $\phi^{(i)}$ is estimated once and reused across the prediction horizon.

Given a choice of $\hat{f}$, the learning problem reduces to estimating the corresponding parameters. Depending on the model class, these parameters may be directly optimized or produced by the mapping $\psi$.

\section{Learning Objective}

All models are trained to minimize trajectory-level prediction error:
\begin{equation}
  \min_{\Theta} \;
  \mathbb{E}_{(s_{t:t+H}, a_{t:t+H-1}, \xi) \sim \mathcal{D}_{\text{train}}}
  \left[
    \sum_{k=1}^{H}
    \mathcal{L}\big(
      s_{t+k},\,
      \hat{s}_{t+k}
    \big)
  \right],
\end{equation}
where $\hat{s}_{t+k}$ is obtained by rolling out the model over the prediction horizon.

The optimization variable $\Theta$ denotes the set of learnable quantities. For fixed-parameter models, $\Theta$ corresponds to model parameters. For context-conditioned models, $\Theta$ denotes the parameters of the mapping $\psi$.

This objective captures compounding prediction errors over the prediction horizon.

\section{Model Structure}

Different modeling approaches correspond to different choices of the function $\hat{f}$ and different assumptions on the structure of $\phi^{(i)}$.

Three categories of models are considered:

\begin{itemize}
\item \textbf{Fixed-parameter models:} $\phi^{(i)} \equiv \phi$ for all $i$, where $\phi$ is a constant parameter vector shared across all trajectory segments.

\item \textbf{Physics-informed models:} $\phi^{(i)}$ represents low-dimensional quantities that encode prior knowledge about system dynamics, corresponding to structured parameterizations.

\item \textbf{Purely data-driven models:} $\phi^{(i)}$ parameterizes high-capacity function approximators with minimal structural constraints, corresponding to flexible parameterizations.
\end{itemize}

This formulation enables a controlled comparison of different model structures under identical data and training conditions.